\section{Restricted Metric Paths in Affine-Invariant Geometry}
\label{sec:spd-rdgc}

We now specialize the path system to positive-definite metrics and make the
terminal optimization problem explicit.  Let
\[
  f(\theta)=\frac12\theta^\top H\theta,
  \qquad H\in\SPD^d,
\]
and let \(G\in\SPD^d\) be a positive metric.  In Hessian-relative
coordinates,
\[
  S=H^{-1/2}GH^{-1/2}.
\]
The Hessian-relative condition ratio is
\[
  \kappa_{\rm gen}(H,G)
  :=\frac{\lambda_{\max}(G^{-1/2}HG^{-1/2})}
  {\lambda_{\min}(G^{-1/2}HG^{-1/2})}
  =\kappa(S).
\]
Whenever the shorthand \(\kappa(G^{-1}H)\) appears, it denotes this ratio of
positive generalized eigenvalues, never the Euclidean singular-value
condition number of the generally nonnormal product \(G^{-1}H\).
The affine-invariant metric and distance are
\[
  \norm{Z}_{S,\AI}
  =\fro{S^{-1/2}ZS^{-1/2}},
  \qquad
  d_{\AI}(S_0,S_1)
  =\fro{\log(S_0^{-1/2}S_1S_0^{-1/2})}.
\]
The cone \(\SPD^d\) is a finite-dimensional Hadamard manifold
\cite{bhatia2007positive,pennec2006riemannian}.

For \(K\ge1\), define the condition target
\[
  \Ck=\{S\in\SPD^d:\kappa(S)\le K\}.
\]
The unrestricted distance to this target has a one-dimensional spectral
formula.

\begin{theorem}[Full-SPD spectral benchmark]
\label{thm:full-spd-benchmark}
Let \(y_1\le\cdots\le y_d\) be the ordered log-eigenvalues of
\(S_0\in\SPD^d\).  Then
\[
  D_K(S_0)
  :=d_{\AI}(S_0,\Ck)
  =
  \min_{c\in\R}
  \left(
  \sum_{i=1}^d\dist(y_i,[c,c+\log K])^2
  \right)^{1/2}.
\]
\end{theorem}

Let \(\F\subset\SPD^d\) be an admissible family and set
\[
  \mathcal S_\F(H)
  =
  \{H^{-1/2}GH^{-1/2}:G\in\F\}.
\]
A realization includes a path-connected component, a treatment of scale
freedom, an admissible absolutely continuous path class, and its induced or
quotient affine-invariant length.

\begin{definition}[Restricted dynamic geometric complexity]
\label{def:restricted-complexity}
For \(G_0\in\F\) and \(S_0=H^{-1/2}G_0H^{-1/2}\), define
\[
\begin{aligned}
  D_{K,\F}(S_0;H)
  =
  \inf\{&\Len_{\AI}(S_{[0,T]}):
  S(0)=S_0,\ S(T)\in\Ck,\\
  &S_t\in\mathcal S_\F(H),
  \ S_{[0,T]}\text{ admissible}\}.
\end{aligned}
\]
The value is \(+\infty\) if the relative target is empty or belongs to a
different admissible component.
\end{definition}

For a closed geodesically convex realization, RDGC is an ordinary Hadamard
projection distance.  This fact connects the nonsmooth length problem to the
fixed-horizon action used in the response theory.

\begin{proposition}[Static RDGC projection]
\label{prop:rdgc-target-projection}
Suppose that \(\mathcal S_\F(H)\) is nonempty, closed, and geodesically
convex and that
\[
  \mathcal T_{K,\F}(H)
  =\mathcal S_\F(H)\cap\Ck
\]
is nonempty.  Then \(S_0\) has a unique projection
\(P_{K,\F}(S_0;H)\) onto \(\mathcal T_{K,\F}(H)\), and
\[
  D_{K,\F}(S_0;H)
  =
  d_{\AI}\bigl(S_0,P_{K,\F}(S_0;H)\bigr).
\]
For quadratic kinetic action with terminal indicator of
\(\mathcal T_{K,\F}(H)\), the value at time \(t<T\) is
\[
  V(t,S)=
  \frac{d_{\AI}(S,\mathcal T_{K,\F}(H))^2}{2(T-t)}.
\]
The minimizing path is the unique constant-speed geodesic to the projection.
\end{proposition}

The projection reduction is a benchmark for unrestricted curves in a closed
geodesically convex realization.  It does not apply when the update protocol
restricts admissible velocities or excludes the ambient geodesic.  The exact
coordinate-sequential model in \cref{prop:sequential-diagonal-rdgc} gives a
structured example whose RDGC strictly exceeds this ambient benchmark.

The convex potentials required by the global response theorem arise
naturally from relative spectra.  For \(H_a\in\SPD^d\) and a convex
permutation-invariant function \(\phi_a:\R^d\to\R\), write
\[
  \Omega_a(G)
  =
  \phi_a\!\left(
  \log\lambda(G^{-1/2}H_aG^{-1/2})
  \right).
\]

\begin{proposition}[Geodesically convex spectral information]
\label{prop:spd-spectral-convexity}
Let \(\F\subset\SPD^d\) be nonempty, closed, and geodesically convex.  If
\(\rho:\R^q\to\R\) is convex and nondecreasing in every coordinate, then
\[
  \mathcal J(G)=\rho(\Omega_1(G),\ldots,\Omega_q(G))
\]
is closed and geodesically convex on \(\F\).  If every
\(\phi_a(x+c\mathbf1)=\phi_a(x)\), then
\(\mathcal J(cG)=\mathcal J(G)\).
In particular,
\[
  \phi_{\rm osc}(x)=\max_i x_i-\min_i x_i
\]
generates \(\mathcal J_H(G)=\log\kappa_{\rm gen}(H,G)\).  Smooth spectral
approximations retain geodesic convexity.
\end{proposition}

Combining this construction with the global path theorem gives the exact SPD
interaction law used below.

\begin{corollary}[SPD--AIRM M\"obius--Jacobi specialization]
\label{cor:spd-mobius-jacobi}
Let \(\mathcal N\) be \(\SPD^d\), or a structured AIRM realization that is
a complete simply connected totally geodesic manifold.  Let the bulk and
terminal potentials in \cref{thm:global-mobius-jacobi} be smooth
geodesically convex spectral information potentials satisfying its
regularity assumptions.  Then the conclusions of that theorem hold with
\[
  J_u\xi
  =-D_t^2\xi-R^{\AI}(\xi,\dot G_u)\dot G_u
  +\operatorname{Hess}^{\AI}W_u\,\xi,
\]
and
\[
  -\ip{R^{\AI}(\xi,\dot G_u)\dot G_u}{\xi}_{G_u}\ge0.
\]
In particular, the pair effect is
\[
  \boxed{
  \zeta_{\{i,j\}}
  =-\int_{[0,1]^2}
  \left\langle
  D\Delta_i,\mathfrak J_u^{-1}D\Delta_j
  \right\rangle\dd u_i\dd u_j.}
\]
Fixed smooth active strata are covered.  The nonsmooth condition-number width
requires a generalized Jacobi inclusion at eigenvalue or active-index
changes.
\end{corollary}

\begin{remark}[Structured families covered directly]
The positive diagonal cone in a fixed basis and every fixed block-diagonal
SPD cone are complete simply connected totally geodesic AIRM submanifolds;
their geodesics, logarithms, and exponentials preserve the declared block
structure.  The determinant-one diagonal family is the corresponding flat
totally geodesic scale section.  Hence the corollaries above apply directly to
these families.  Generic low-rank and Kronecker parameterizations require
their own embedded or quotient geometry and are outside this direct
totally-geodesic specialization.
\end{remark}

The hard condition target itself also has a classical response away from
spectral collisions.  It is useful here to write
\[
  \kappa_{\rm gen}(H,G)
  =\frac{\lambda_{\max}(G^{-1/2}HG^{-1/2})}
  {\lambda_{\min}(G^{-1/2}HG^{-1/2})}.
\]

\begin{corollary}[Moving hard-condition target on a smooth spectral stratum]
\label{cor:hard-condition-response}
Let \(\F\) be a complete simply connected totally geodesic AIRM
realization, fix \(G_0\in\F\), and let \(H_u\in\SPD^d\) and \(K_u>1\) vary
smoothly.  Define
\[
  c_u(G)=\log\kappa_{\rm gen}(H_u,G)-\log K_u,
  \qquad
  \mathcal T_u=\{G\in\F:c_u(G)\le0\}.
\]
Suppose that \(G_0\notin\mathcal T_{u_0}\), that the largest and smallest
generalized eigenvalues of \((H_{u_0},q_0)\) are simple at the AIRM projection
\(q_0=P_{\mathcal T_{u_0}}G_0\), and that
\(\dd(c_{u_0}|_{\F})(q_0)\ne0\).  Then, locally in \(u\), the projection
endpoint \(q_u\), its constant-speed geodesic \(\gamma_u\), and its positive
terminal multiplier are smooth.  Moreover,
\[
  \mathcal V_{\rm hard}(u)
  =\frac{d_{\AI}(G_0,\mathcal T_u)^2}{2T}
\]
satisfies \textup{(HT4)}--\textup{(HT5)} with
\(\mathfrak J_u\) equal to the kinetic index form plus
\(\lambda_u\operatorname{Hess}^{\AI}(c_u|_{\F})\) at the endpoint.
If the active defining function is affine in the intervention coordinates on
a Boolean face, and the same active spectral stratum, simple extreme
eigenvalues, LICQ, positive multiplier, and bordered invertibility persist on
that entire face, \textup{(HT6)} gives its exact hard-RDGC interaction.

Thus the response differentiates the moving hard-target projection itself.
The conclusion ends precisely when an extreme generalized eigenvalue loses
simplicity, the restricted boundary loses regularity, strict complementarity
fails, or the active stratum changes.
\end{corollary}

